% we include the glossary here (frontmatter is included with \input, so this command is as if it was in main.tex)
%All acronyms must be written in this file.

% Institutional
\newacronym{ISEP}{ISEP}{Instituto Superior de Engenharia do Porto}
\newacronym{UC}{UC}{Unidade Curricular}
\newacronym{PESTI}{PESTI}{Projeto/Estágio}
\newacronym{LEI-PROJ}{LEI-PROJ}{Projeto em Licenciatura em Engenharia Informática}

% Cloud & Infrastructure
\newacronym{AWS}{AWS}{Amazon Web Services}
\newacronym{CI/CD}{CI/CD}{Continuous Integration / Continuous Deployment}
\newacronym{SES}{SES}{Simple Email Service}
\newacronym{WAF}{WAF}{Web Application Firewall}
\newacronym{JWKS}{JWKS}{JSON Web Key Set}

% Development & Architecture
\newacronym{API}{API}{Application Programming Interface}
\newacronym{REST}{REST}{Representational State Transfer}
\newacronym{CRUD}{CRUD}{Create, Read, Update, Delete}
\newacronym{E2E}{E2E}{End to End}
\newacronym{FURPS+}{FURPS+}{Functionality, Usability, Reliability, Performance, Supportability}
\newacronym{MVP}{MVP}{Minimum Viable Product}
\newacronym{UI}{UI}{User Interface}
\newacronym{UX}{UX}{User Experience}
\newacronym{QA}{QA}{Quality Assurance}
\newacronym{IA}{IA}{Inteligência Artificial}

% Database
\newacronym{PostgreSQL}{PostgreSQL}{PostgreSQL Relational Database}

% Security & Identity
\newacronym{KYC}{KYC}{Know Your Customer}
\newacronym{KYB}{KYB}{Know Your Business}
\newacronym{MFA}{MFA}{Multi-Factor Authentication}
\newacronym{RBAC}{RBAC}{Role-Based Access Control}
\newacronym{TTL}{TTL}{Time To Live}

% Metrics
\newacronym{KPI}{KPI}{Key Performance Indicator}

% Protocols & Diagrams
\newacronym{HTTP}{HTTP}{Hypertext Transfer Protocol}
\newacronym{SD}{SD}{Sequence Diagram}
\newacronym{SSD}{SSD}{System Sequence Diagram}
\newacronym{SSE}{SSE}{Server-Sent Events}

% Other
\newacronym{CSV}{CSV}{Comma-Separated Values}
\newacronym{OJC}{OJC}{Only Just Causes}
\newacronym{PDF}{PDF}{Portable Document Format}


\frontmatter % Use roman page numbering style (i, ii, iii, iv...) for the pre-content pages

\pagestyle{plain} % Default to the plain heading style until the thesis style is called for the body content

%----------------------------------------------------------------------------------------
%	TITLE PAGE
%----------------------------------------------------------------------------------------

\maketitlepage


%----------------------------------------------------------------------------------------
%	STATEMENT of INTEGRITY
%----------------------------------------------------------------------------------------
\integritystatement

%----------------------------------------------------------------------------------------
%	DEDICATION  (optional)
%----------------------------------------------------------------------------------------
%
%\dedicatory{For/Dedicated to/To my\ldots}
%\begin{dedicatory}
%\end{dedicatory}

%----------------------------------------------------------------------------------------
%	ABSTRACT PAGE
%----------------------------------------------------------------------------------------

\begin{abstract}

Qualquer aplicação que pretende ser utilizada em larga escala, por diferentes tipos de utilizadores e para diferentes propósitos, necessita de uma camada de operação interna que permita às suas equipas gerir o que acontece por detrás do produto. A OnlyHighIQ desenvolve um ecossistema de aplicações cujo primeiro produto é a JustCauses, uma plataforma de \textit{crowdfunding} para apoio a causas sociais e, para suportar a operação deste ecossistema de forma estruturada, tornou-se necessário conceber uma solução de \textit{backoffice} centralizada. É neste contexto que se enquadra o trabalho descrito neste relatório, desenvolvido durante um estágio na OnlyHighIQ no âmbito da Licenciatura em Engenharia Informática do ISEP.

O objetivo do estágio foi conceber e implementar um \textit{backoffice} multiaplicação para o ecossistema da empresa, tendo a JustCauses sido a primeira aplicação integrada. A solução permite às equipas de operação gerir campanhas, categorias, verificações de identidade e verificações de organizações, integrando a autenticação centralizada existente com um modelo interno de autorização baseado em perfis, permissões por aplicação e registo de auditoria. A sua arquitetura modular foi pensada para suportar futuras aplicações do ecossistema sem alterações estruturais de fundo.

Ao longo das \textit{sprints}, o trabalho foi dividido entre a base transversal do \textit{backoffice} e os módulos aplicados à JustCauses. Na base comum ficaram a integração com a autenticação AWS Cognito, a gestão de operadores administrativos, a atribuição de perfis e permissões por aplicação, a configuração de acesso por página e a auditoria de atividade. Na JustCauses foram tratados fluxos como revisão de campanhas com comunicação por e-mail, gestão de categorias, verificação \gls{KYC}/\gls{KYB} e indicadores operacionais. A solução usa React no \textit{front-end} e Node.js/Express com TypeScript no \textit{back-end}, com persistência em PostgreSQL e DynamoDB e integração com serviços AWS como Cognito, S3, SES e CloudWatch.

No final do período coberto por este relatório, a solução encontrava-se operacional no âmbito definido para o estágio e integrada no processo de validação da empresa. O estágio permitiu também consolidar competências de desenvolvimento \textit{full-stack}, arquitetura de \textit{software}, integração com serviços \textit{cloud} e entrega iterativa em contexto real de produto.

\vspace{20mm}

\textbf{Palavras-chave (Tema)}: Backoffice, Multiaplicação, Autorização, Auditoria.

\textbf{Palavras-chave (Tecnologias)}: React, Node.js, TypeScript, AWS, PostgreSQL, DynamoDB.

\end{abstract}

\begin{abstractotherlanguage}

Any application that aims to be used at scale, by different types of users and for different purposes, requires an internal operations layer that allows its teams to manage what happens behind the product. OnlyHighIQ is building an ecosystem of applications whose first product is JustCauses, a crowdfunding platform for social causes, and to support the operation of this ecosystem in a structured way, a centralised backoffice solution became necessary. This report describes the work carried out during an internship at OnlyHighIQ as part of the BSc in Informatics Engineering at ISEP.

The goal of the internship was to design and implement a multi-application backoffice for the company's ecosystem, with JustCauses as the first integrated application. The solution allows operations teams to manage campaigns, categories, identity checks, and organization verification, integrating the existing centralised authentication with an internal authorisation model based on roles, per-application permissions, and audit logging. Its modular architecture was designed to accommodate future applications in the ecosystem without major structural changes.

Across the sprints, the work moved between the common backoffice foundation and the modules applied to JustCauses. The shared foundation included integration with AWS Cognito authentication, administrative operator management, per-application roles and permissions, page access configuration, and activity auditing. The JustCauses work covered campaign review with email communication, category management, KYC/KYB verification, and operational indicators. The solution uses React on the front end and a Node.js/Express back end written in TypeScript, with PostgreSQL, DynamoDB, and AWS services such as Cognito, S3, SES, and CloudWatch.

By the end of the internship period, the backoffice was operational within the defined scope and integrated into the company's validation process. The internship also strengthened practical skills in full-stack development, software architecture, cloud service integration, and iterative delivery in a real product environment.

\vspace{20mm}

\textbf{Keywords (Theme)}: Backoffice, Multi-application, Authorization, Auditing.

\textbf{Keywords (Technologies)}: React, Node.js, TypeScript, AWS, PostgreSQL, DynamoDB.

\end{abstractotherlanguage}


%----------------------------------------------------------------------------------------
%	ACKNOWLEDGEMENTS (optional)
%----------------------------------------------------------------------------------------

\begin{acknowledgements}

A realização deste estágio e deste relatório beneficiou do acompanhamento e da disponibilidade de várias pessoas, a quem deixo um agradecimento sincero.

Agradeço ao Professor Paulo Maio pelo acompanhamento académico e pelas observações feitas ao longo do trabalho.

Agradeço também a toda a equipa da OnlyHighIQ pela integração num contexto real de desenvolvimento, pela partilha de conhecimento e pela confiança dada durante a implementação do backoffice da JustCauses.

Aos colegas com quem trabalhei durante o estágio, agradeço o apoio diário, a entreajuda nas fases de análise, implementação e validação, e o contexto técnico e funcional que tornou possível avançar com o projeto de forma consistente.

Aos amigos que conheci no curso, que fizeram destes três anos de estudos, praxes e saídas à noite a melhor experiência que um estudante poderia ter. Agora que olho para trás, só me arrependo de não ter aproveitado ainda mais.

Agradeço também a toda a minha família, em especial ao meu pai, também ele licenciado neste mesmo curso, que de certa forma me transmitiu a paixão por este mundo. Agradeço ainda, como é obvio, aos meus gatos, a Simone e o Becas.

Por fim, agradeço à minha namorada, que me acompanhou neste último ano. Sei que, graças a ela, fui mais responsável e consegui manter o foco necessário para concluir este trabalho. Amo-te, Bea, e espero que, para o ano, possa ler a tua dedicatória no teu relatório de estágio.

\end{acknowledgements}


%----------------------------------------------------------------------------------------
%	LIST OF CONTENTS/FIGURES/TABLES PAGES
%----------------------------------------------------------------------------------------

\tableofcontents % Prints the main table of contents

\listoffigures % Prints the list of figures

\listoftables % Prints the list of tables

\iflanguage{portuguese}{
	\renewcommand{\listalgorithmname}{Lista de Algor\'itmos}
}
\listofalgorithms % Prints the list of algorithms
%\addchaptertocentry{\listalgorithmname}


\renewcommand{\lstlistlistingname}{List of Source Code}
\iflanguage{portuguese}{
	\renewcommand{\lstlistlistingname}{Lista de C\'odigo}
}
\lstlistoflistings % Prints the list of listings (programming language source code)
%\addchaptertocentry{\lstlistlistingname}


%----------------------------------------------------------------------------------------
%	ABBREVIATIONS
%----------------------------------------------------------------------------------------
%\begin{abbreviations}{ll} % Include a list of abbreviations (a table of two columns)
%%\textbf{LAH} & \textbf{L}ist \textbf{A}bbreviations \textbf{H}ere\\
%%\textbf{WSF} & \textbf{W}hat (it) \textbf{S}tands \textbf{F}or\\
%\end{abbreviations}

%----------------------------------------------------------------------------------------
%	SYMBOLS
%----------------------------------------------------------------------------------------

%\begin{symbols}{lll} % Include a list of Symbols (a three column table)

%$a$ & distance & \si{\meter} \\
%$P$ & power & \si{\watt} (\si{\joule\per\second}) \\
%%Symbol & Name & Unit \\

%\addlinespace % Gap to separate the Roman symbols from the Greek

%$\omega$ & angular frequency & \si{\radian} \\

%\end{symbols}

%----------------------------------------------------------------------------------------
%	ACRONYMS
%----------------------------------------------------------------------------------------

\newcommand{\listacronymname}{List of Acronyms}
\iflanguage{portuguese}{
	\renewcommand{\listacronymname}{Lista de Acr\'onimos}
}

%Use GLS
\glsresetall
\printglossary[title=\listacronymname,type=\acronymtype,style=long]

%----------------------------------------------------------------------------------------
%	DONE
%----------------------------------------------------------------------------------------

\mainmatter % Begin numeric (1,2,3...) page numbering
\pagestyle{thesis} % Return the page headers back to the "thesis" style
